% Part: first-order-logic
% Chapter: model-theory
% Section: overspill

\documentclass[../../../include/open-logic-section]{subfiles}

\begin{document}

\olfileid{mod}{bas}{ove}
\olsection{Overspill}

\begin{thm}
\ollabel{overspill} Yen himpunan ukara $\Gamma$
nduweni model winates kang ukurane bisa
ngluwihi sembarang wates, mula himpunan
iku nduweni model tanpa wates.
\end{thm}

\begin{proof}
Jembarake basa kanggo $\Gamma$ kanthi
nambah konstanta anyar kang cacahé
kaetung, yaiku $c_0$, $c_1$, \dots
lan gatekna himpunan $\Gamma \cup \{c_i \neq c_j :
i \neq j\}$. Pratelan yen $\Gamma$
nduweni model winates kang ukurane
bisa ngluwihi sembarang wates
ateges kanggo saben $m >0$
ana $n\ge m$ saéngga $\Gamma$
nduweni model kang kardinalitase~$n$.
Iki nuduhake yen saben subhimpunan
winates saka $\Gamma \cup \{c_i \neq
c_j : i \neq j\}$ bisa dipenuhi.
Miturut kompaktitas, $\Gamma
\cup \{c_i \neq c_j : i \neq j\}$
nduweni model $\Struct M$
kang ranahe mesthi tanpa wates,
amarga model iku nyukupi
kabeh pertidaksamaan $c_i \neq c_j$.
\end{proof}

\begin{prop}
\ollabel{inf-not-fo}
Ora ana ukara $!A$ ing sembarang
basa orde kapisan kang bener ing
!!a{structure}~$\Struct M$
yen lan mung yen ranah
$\Domain{M}$ saka
!!{structure} iku tanpa wates.
\end{prop}

\begin{proof}
Yen ana ukara $!A$ kaya mangkono,
negasine $\lnot !A$ mesthi bener
ing kabeh !!{structure}s winates
lan mung ing struktur-struktur
iku. Mula negasi kasebut
nduweni model winates kang
ukurane bisa ngluwihi
sembarang wates, nanging
ora nduweni model tanpa
wates. Iki mbantah
\olref{overspill}.
\end{proof}

\end{document}
